% Generated draft: requires mathematical and source review.
\documentclass[11pt]{article}
\usepackage[margin=25mm]{geometry}
\usepackage{fontspec}
\setmainfont{Latin Modern Roman}
\usepackage{amsmath,amssymb,mathtools,unicode-math}
\setmathfont{Latin Modern Math}
\usepackage{enumitem,needspace,xurl,hyperref}
\setlength{\parindent}{0pt}
\setlength{\parskip}{5pt}
\setlength{\emergencystretch}{3em}
\newcommand{\N}{\mathbb{N}}
\newcommand{\Z}{\mathbb{Z}}
\newcommand{\Q}{\mathbb{Q}}
\newcommand{\R}{\mathbb{R}}
\newcommand{\C}{\mathbb{C}}
\newcommand{\F}{\mathbb{F}}
\newcommand{\A}{\mathbb{A}}
\newcommand{\PP}{\mathbb{P}}
\newcommand{\E}{\mathbb{E}}
\newcommand{\eps}{\varepsilon}
\newcommand{\dd}{\,\mathrm{d}}
\begin{document}
\section*{UnsolvedMath ID: 30000589}
\textbf{Source problem code:} catalogue entry
\subsection{Definitions and mathematical statement}
\textbf{Two chain-level questions.} For the construction of homotopy orbits of an \(S^1\)-action, do analogues of the source's primitive and cobar quasi-isomorphism results hold for every compact Lie group? What are the mathematical meaning and content of the resulting higher operations \(\omega_k\)? The specific circle-action construction, comparison maps, and operations are not reproduced, so a general homotopy-orbit construction alone is not substituted for the requested stronger analogues.\par For a compact Lie group \(G\) acting on a space \(Y\), its homotopy-orbit space is conventionally \(Y_{hG}=EG\times_GY\), with \(EG\) contractible and free. It is not identified with the ordinary orbit space \(Y/G\) without the appropriate hypotheses. With a specified coefficient ring \(R\), chains \(C_*(Y_{hG};R)\) provide the target of a chain-level comparison. A quasi-isomorphism is a chain map inducing an isomorphism in every homology degree; the exact source complexes and maps must be stated to formulate the desired analogues.\par In a coaugmented differential graded coalgebra \(C\), an element is primitive when its coproduct is \(1\otimes c+c\otimes1\). A standard cobar construction is the tensor algebra on the desuspension of the reduced coalgebra, \(\Omega C=T(s^{-1}\overline C)\), with differential determined by the internal differential and reduced coproduct. The source does not identify which coalgebra, primitive construction, completeness, connectivity, or convergence conventions its results use. Integral and rational settings, connected and disconnected compact Lie groups, and possible torsion hypotheses are not silently equated.\par The symbols \(\omega_k\) have no supplied domains, codomains, arities, degrees, identities, or dependence on choices. Explaining their mathematical content requires those definitions and the intended comparison with geometric or homological invariants. They are not arbitrarily called Massey products, brackets, or operations of a chosen infinity structure. Both the all-compact-Lie-group generalization and the higher-operation interpretation remain explicit tasks pending recovery of the circle-case statements.\par The supplied triage checked 2026-08-21 labels the item open and reports constructions in special settings without a located general resolution. The later context reference contains no quoted theorem in this record. This dated assessment requires independent verification.
\subsection{Short English statement}
Generalize the specific primitive and cobar comparison results from circle actions to compact Lie groups, and explain the resulting higher operations. The source's complexes, coefficients, comparison maps, and operation definitions are missing; the ordinary existence of a homotopy-orbit space does not answer those questions.
\subsection{Sources}
\begin{itemize}\item Local UnsolvedMath registry, permanent ID 30000589, code catalogue entry; \emph{Homotopy-Orbit Constructions for Compact Lie-Group Actions}, Oberwolfach Reports Open Problems.\item \emph{Topologie (2007)}, Oberwolfach Reports, DOI 10.4171/owr/2006/43, pp. 2579--2630. \url{https://doi.org/10.4171/owr/2006/43}. Proposed year and year tag: 2006; no proposer named.\item Supplied literature label \emph{original OWR report}: \url{https://doi.org/10.4171/owr/2006/43}.\item Supplied literature label \emph{later higher-operation context}: \url{https://doi.org/10.1007/s00208-023-02741-4}.\item Literature triage checked 2026-08-21, status \texttt{open}; record updated 2026-09-16. The \texttt{corrected\_verified} field records notation normalization and preservation of both questions. Verify dated claims independently.\end{itemize}
\end{document}
